% =====================================================================================================
% PART 5 (frames 19-23): the barrier, how the results were obtained, open problems, summary.
% Sources in the paper: Section 1.2 (paragraphs "A barrier", "How these results were obtained"),
% Section 7 (Theorem 7.1, the distribution mu*, Lemma 7.2, Table 2, Figure 5), Section 8,
% Appendix B and the AI methodology statement.
% Figures: figs/p5-barrier.tex
% Citation tags are the alpha labels of the paper's bibliography: [TSSW00], [HHM+17], [Cha16], [AOTW14].
% =====================================================================================================

% local helpers of part 5 (names are prefixed with pV to avoid clashes with other parts)
\providecommand{\pVnum}[1]{\tikz[baseline=(n.base)]\node[circle,fill=acc,text=white,font=\bfseries\large,
  inner sep=0pt,minimum size=8.5mm] (n) {#1};}
\providecommand{\pVto}{\ensuremath{\longrightarrow}}
\providecommand{\pVimp}{{\color{acc}$\boldsymbol{\Longrightarrow}$}}
% small swatch of a cell of the figure of frame 20; #1 = fill colour
\providecommand{\pVswatch}[1]{\tikz[baseline=-0.6ex]\draw[draw=black!45,line width=0.5pt,fill=#1]
  (-2.2mm,-1.3mm) rectangle (2.2mm,1.3mm);}

% ---- Frame 19: A barrier at 27/29
\begin{frame}{A barrier at $27/29$}
\vspace{-0.5mm}
\begin{block}{Theorem 7.1\ \ (computer-assisted)}
For \textbf{every} family $(I_0,I_1,I_2)$ of \MC{3} gadgets ---
any number of auxiliary vertices,\\ with or without global vertices:
\par\vspace{0.6mm}
\hspace*{5mm}%
\mbox{\Large$\displaystyle\max_{x\in\Z_3^3}\,S_{\yc}(x)\;\;\ge\;\;
  {\color{acc}\boldsymbol{\frac{27}{29}}}\,\sum_{t}C_t$}%
\hfill
\begin{minipage}[c]{71mm}\small\color{soft}\raggedright
$S_{\yc}(x)=\sum_tV_t(x,\yc)$:\ \ all three gadgets\par\vspace{0.2mm}
on one colour pattern $x$, with canonical globals\par\vspace{1mm}
$C_t$:\ \ completeness of the type-$t$ gadget\par\vspace{0.2mm}
\hphantom{$C_t$:\ \ }(its least value on satisfying colourings)
\end{minipage}
\vspace{-0.6mm}
\end{block}
\vspace{2.6mm}
\begin{minipage}[c]{11mm}\pVimp\end{minipage}%
\begin{minipage}[c]{130mm}\large\raggedright
For \MC{3}, the criterion cannot give a factor below $27/29$:\par\vspace{0.5mm}
\hl{the simple gadgets $E$ and $N$ are optimal}
\end{minipage}
\par\vspace{3.1mm}
\begin{minipage}[c]{11mm}\pVimp\end{minipage}%
\begin{minipage}[c]{130mm}\large\raggedright
With tightness \paperref{Proposition 4.3}: $27/29$ is \hl{the limit for gadget reductions}\par\vspace{0.5mm}
from H\aa stad's equations that use his theorem as a black box
\end{minipage}
\par\vspace{3.1mm}
\hspace*{11mm}{\color{soft}Proof: LP duality, in the spirit of} \src{TSSW00, HHM$^{+}$17}
\end{frame}

% ---- Frame 20: The proof: one distribution that is bad for every gadget
\begin{frame}{The proof: one distribution that is bad for every gadget}
\vspace{-0.3mm}
\centerline{% Frame 20: the distribution mu* of the barrier (Theorem 7.1) on the classes A, D, T_1, T_2.
% Adapted from ../paper/figures/barrier.tex (Figure 5 of the paper).  Mathematical content unchanged:
% same boxes (x-units as in the paper: A=[0,1], D=[1.75,3.75], T_1=[4.5,7.5], T_2=[8.25,11.25]), one dot
% per tuple at the same positions with area proportional to its mu*-mass (5/87 on A, 1/29 = 3/87
% elsewhere, radii in the ratio sqrt(5/3) as in the paper), same class masses, same cells (sat_t shaded;
% conditional probabilities 9/18, 9/18 and 5/20, 6/20, 9/20), same right-hand column.
% Changes for the slide (layout only): larger fonts (the fractions in the cells are scaled by 1.15),
% x-unit 0.82cm, lower boxes and taller cells; the mass label of the narrow box A is moved 2.4mm to the
% left; in the right-hand column the certificate's guarantee (8/9, 9/10) is set in the accent colour and
% the result 27/29 in bold; the description of mu* (Section 7 of the paper: uniform within each class,
% class masses 5:6:9:9) is added at the top right.
% The legend of the cells (caption of Figure 5) is given by the three text zones under the figure on
% the slide.  The data lines of the \foreach lists below were compared with the paper's file by a script.
\begingroup
\def\pVrh{0.32}\def\pVbt{1.16}% row height of the dots; top of the boxes (3*0.32 + 2*0.1)
\def\pVra{1.3mm}%      dot radius on A (mass 5/87)
\def\pVrb{1.00698mm}% dot radius elsewhere (mass 1/29 = 3/87) = 1.3mm*sqrt(3/5); area ratio 5:3
\def\pVch{0.6}\def\pVcp{0.69}% height and pitch of the cells
\begin{tikzpicture}[x=0.82cm,y=1cm,
  cls/.style={draw,rounded corners=2.5pt,line width=0.7pt,fill=black!3},
  cname/.style={font=\Large,inner sep=1pt,text height=1.7ex,text depth=0.5ex},
  cmass/.style={font=\normalsize,inner sep=1pt},
  rlab/.style={font=\normalsize,anchor=east,inner sep=1pt},
  cell/.style={draw=black!45,line width=0.5pt},
  hi/.style={cell,fill=black!12},
  val/.style={font=\normalsize,scale=1.15},
  satv/.style={font=\normalsize},
  bnd/.style={font=\normalsize,anchor=west,inner sep=1pt,scale=1.12}]
% ---------------------------------------------------------------- boxes and dots
\foreach \l/\r in {0/1, 1.75/3.75, 4.5/7.5, 8.25/11.25}
  \draw[cls] (\l,0) rectangle (\r,\pVbt);
\foreach \i in {0,1,2}{
  \fill[black!60] (0.5,\pVbt-0.1-\i*\pVrh-0.5*\pVrh) circle[radius=\pVra];
  \foreach \x in {2.25,3.25,5,6,7,8.75,9.75,10.75}
    \fill[black!60] (\x,\pVbt-0.1-\i*\pVrh-0.5*\pVrh) circle[radius=\pVrb];}
\node[rlab] at (-0.25,0.5*\pVbt) {$\mu^\ast$};
% ---------------------------------------------------------------- headers
\foreach \c/\n/\m in {0.5/{\mathcal A}/{\tfrac{5}{29}=3\cdot\tfrac{5}{87}},
                      2.75/{\mathcal D}/{\tfrac{6}{29}=6\cdot\tfrac{1}{29}},
                      6/{\mathcal T_1}/{\tfrac{9}{29}=9\cdot\tfrac{1}{29}},
                      9.75/{\mathcal T_2}/{\tfrac{9}{29}=9\cdot\tfrac{1}{29}}}{
  \node[cname] at (\c,\pVbt+0.75) {$\n$};
  \ifdim\c pt<1pt% the label of the narrow box A is moved 2.4mm to the left, away from the label of D
    \node[cmass,xshift=-2.4mm] at (\c,\pVbt+0.29) {$\m$};
  \else
    \node[cmass] at (\c,\pVbt+0.29) {$\m$};
  \fi}
% ---------------------------------------------------------------- description of mu* (top right)
\node[anchor=north west,font=\normalsize,inner sep=1pt,align=left] at (11.55,\pVbt+0.96)
  {$\mu^\ast$: \textbf{uniform} within\\[0.3mm]each class, masses\\[1.1mm]\hspace*{4mm}{\large$\mathbf{5:6:9:9}$}};
% ---------------------------------------------------------------- certificate grid
\node[font=\normalsize,anchor=south west,inner sep=1pt] at (11.55,-0.17) {$\mathbb{E}_{\mu^\ast}V_t\,/\,C_t\ \ge$};
\foreach \t/\a/\d/\p/\q/\bd in {
   0/sat/sat/{\tfrac{9}{18}}/{\tfrac{9}{18}}/{\tfrac{11}{29}+\tfrac{18}{29}\cdot{\color{acc}\boldsymbol{\tfrac{8}{9}}}=\boldsymbol{\tfrac{27}{29}}},
   1/{\tfrac{5}{20}}/{\tfrac{6}{20}}/sat/{\tfrac{9}{20}}/{\tfrac{9}{29}+\tfrac{20}{29}\cdot{\color{acc}\boldsymbol{\tfrac{9}{10}}}=\boldsymbol{\tfrac{27}{29}}},
   2/{\tfrac{5}{20}}/{\tfrac{6}{20}}/{\tfrac{9}{20}}/sat/{\tfrac{9}{29}+\tfrac{20}{29}\cdot{\color{acc}\boldsymbol{\tfrac{9}{10}}}=\boldsymbol{\tfrac{27}{29}}}}{
  \pgfmathsetmacro\yt{-0.23-\t*\pVcp}
  \pgfmathsetmacro\yb{\yt-\pVch}
  \node[rlab] at (-0.25,{(\yt+\yb)/2}) {$t=\t$};
  \foreach \l/\r/\v in {0/1/\a, 1.75/3.75/\d, 4.5/7.5/\p, 8.25/11.25/\q}{
    \edef\tmp{\v}\def\satstr{sat}%
    \ifx\tmp\satstr
      \draw[hi] (\l,\yt) rectangle node[satv] {$\sat_{\t}$} (\r,\yb);
    \else
      \draw[cell] (\l,\yt) rectangle node[val] {$\v$} (\r,\yb);
    \fi}
  \node[bnd] at (11.55,{(\yt+\yb)/2}) {$\bd$};}
\end{tikzpicture}
\endgroup
}
\par\vspace{1.3mm}
\begin{columns}[T,onlytextwidth]
\begin{column}{0.2\textwidth}
\pVswatch{black!12}\ \textbf{Shaded cells}\par\vspace{1.2mm}
$V_t\ge C_t$ on $\sat_t$,\\ by definition
\end{column}
\begin{column}{0.4pt}\color{acc!25}\rule{0.4pt}{19.5mm}\end{column}
\begin{column}{0.408\textwidth}
\pVswatch{white}\ \textbf{Other cells}\ \ ($\mu^\ast$ given $x\notin\sat_t$)\par\vspace{1.2mm}
an \hl{LP-dual certificate} guarantees,\\ on average,
${\color{acc}\boldsymbol{\tfrac89}}\,C_0$ (type $0$) and\\
${\color{acc}\boldsymbol{\tfrac9{10}}}\,C_t$ (types $1,2$)\hfill\paperref{Lemma 7.2}
\end{column}
\begin{column}{0.4pt}\color{acc!25}\rule{0.4pt}{19.5mm}\end{column}
\begin{column}{0.345\textwidth}
\textbf{Exact rational certificates}\par\vspace{1.2mm}
$78$, $162$, $162$ rows;\ \ verified in\\ exact arithmetic by\\
\hl{two independent programs}
\end{column}
\end{columns}
\end{frame}

% ---- Frame 21: How these results were obtained
\begin{frame}{How these results were obtained}
\vspace{-1mm}
This research was carried out with \hl{extensive use of an AI system} (Claude, Anthropic),\\
under the direction of the author.
\hfill{\small\color{soft}Details: Appendix B and the AI methodology statement}
\par\vspace{1.2mm}
\setlength{\leftmargini}{5mm}%
\begin{columns}[T,onlytextwidth]
\begin{column}{0.345\textwidth}
\begin{block}{The AI system}
\begin{itemize}\setlength{\itemsep}{4.2mm}
\item proposed the main ideas
\item wrote and ran the search and verification programs
\item searched the literature
\item drafted the paper
\end{itemize}
\end{block}
\end{column}
\begin{column}{0.61\textwidth}
\begin{block}{How the claims were checked}
\begin{itemize}\setlength{\itemsep}{1.7mm}
\item \hl{Independent AI agents} tried to refute every main claim:
  {\small\color{soft}re-read the primary sources, re-derived the proofs,
  re-implemented the computations}
\item Every gadget value and certificate: verified by\\
  \hl{$\ge2$ independent programs}, in exact arithmetic
\item Theorems 1.1--1.3 have \hl{short proofs by hand};\\ only Theorem 7.1 is computer-assisted
\end{itemize}
\end{block}
\end{column}
\end{columns}
\par\vspace{4.6mm}
\vspace{-2.4mm}\hl{The author verified all content} and is responsible for it.\par\vspace{0.4mm}
Code, gadget families, certificates:\ \
\href{https://github.com/ziangni-sys/max-3-cut}{\color{acc}\texttt{github.com/ziangni-sys/max-3-cut}}
\end{frame}

% ---- Frame 22: Open problems
\begin{frame}{Open problems}
\vspace{0.6mm}
\renewcommand{\arraystretch}{1.0}
\begin{tabular}{@{}m{12mm}@{\hspace{2mm}}m{127mm}@{}}
\pVnum{1} & {\large\textbf{Perfect completeness.}\ \ Beat $16/17$ \src{AOTW14} on $3$-colourable graphs?}\\[0.3mm]
          & {\small\color{soft}combine the negation trick with a source that has perfect completeness}\\[2.5mm]
\pVnum{2} & {\large\textbf{Other sources.}\ \ Below $27/29$: change the source, reduction, or analysis}\\[0.3mm]
          & {\small\color{soft}candidates: Chan's predicates \src{Cha16}, where the negation trick applies as well}\\[2.5mm]
\pVnum{3} & {\large\textbf{$\boldsymbol{k\ge4}$.}\ \ Are $45/47$ and $33/34$ optimal for the criterion?}\\[0.3mm]
          & {\small\color{soft}so far, no barrier certificate}\\[2.5mm]
\pVnum{4} & {\large\textbf{Large $\boldsymbol k$.}\ \ Can the criterion improve the constant in $1-\Theta(1/k)$?}\\[0.3mm]
          & {\small\color{soft}the family of this paper gives only $1-\Theta(1/k^2)$}
\end{tabular}
\par\vspace{3.2mm}
{\small\color{soft}\dots\ and $k=3$:}\ \ $0.836$
{\small\color{soft}(optimal under the Unique Games Conjecture)}\ \ vs.\ \ $0.931$ {\small\color{soft}(NP-hard)}
\end{frame}

% ---- Frame 23: Summary
\begin{frame}{Summary}
\vspace{3mm}
\renewcommand{\arraystretch}{1.0}
\begin{tabular}{@{}m{6mm}@{}m{135mm}@{}}
{\color{acc}\rule{1mm}{11.5mm}} & {\Large \MC{3} is NP-hard to approximate within
   \hl{$\mathbf{27/29}+\eps\approx0.931$}}\par\vspace{1.4mm}
   {\color{soft}was $16/17\approx0.941$ since 2012 \src{AOTW14}}\\[8mm]
{\color{acc}\rule{1mm}{11.5mm}} & {\Large \hl{Negation trick}: pair every assignment $X$ with $-X$}\par\vspace{1.4mm}
   {\color{soft}\pVto\ \ colour-symmetric predicates, tiny gadgets}\\[8mm]
{\color{acc}\rule{1mm}{11.5mm}} & {\Large A tight \hl{soundness criterion}, and a matching \hl{barrier} at $27/29$}\par\vspace{1.4mm}
   {\color{soft}\pVto\ \ more colours; with Theorem 1.3: better bounds for every $k$ from $3$ to $20$}
\end{tabular}
\par\vspace{9mm}
{\LARGE\color{acc}\textbf{Thank you!}}\hfill
{\color{soft}code and certificates:}\ \
\href{https://github.com/ziangni-sys/max-3-cut}{\color{acc}\texttt{github.com/ziangni-sys/max-3-cut}}
\end{frame}
